\section{Síntese dos Resultados}
\label{sec:sintese-resultados}

Em relação à QE-1: em que medida a forma de escrita dos cenários BDD influencia a eficácia e a qualidade do código de testes automatizados gerado por IA no contexto do CAPJu? Conforme definido na Seção~\ref{sec:questao-pesquisa}, a eficácia é medida pela completude estrutural do código gerado (Métricas 1.1 a 1.8) e a qualidade, pela ocorrência de erros sintáticos e de execução (Métrica 1.9), ambas adaptadas de \citeonline{Nettur2025}.

Dos 37 cenários BDD planejados, 27 puderam ser executados contra o CAPJu, e os outros 10 foram bloqueados por \textit{bugs} reais da aplicação (4) ou por CAs ambíguos (6). Nos 27 cenários executados, a IA gerou código estruturalmente completo em todas as tentativas, mas apenas 12 (44,4\%) executaram sem erros na primeira execução. Em mediana, foram necessárias duas tentativas por cenário, em uma distribuição assimétrica e não normal, com máximo de seis. A Tabela~\ref{tab:sintese} consolida os principais indicadores.

\begin{table}[htb]
\centering
\caption{Síntese dos principais resultados}
\label{tab:sintese}
\begin{tabular}{p{5.0cm}p{7.6cm}c}
\hline
\textbf{Indicador} & \textbf{Resultado} & \textbf{Seção} \\
\hline
Cenários executados e bloqueados & 27 de 37 executados (73,0\%); 10 bloqueados (27,0\%) & \ref{sec:consolidacao-resultados} \\
Tentativas até o sucesso & Mediana 2 (Q1--Q3: 1--3); máximo 6; distribuição não normal & \ref{sec:analise-descritiva} \\
Correções que alteraram o BDD & 13 de 15 cenários (86,7\%) & \ref{subsec:correcoes} \\
Falhas ligadas ao elemento de interface & 17 de 27 (63,0\%) & \ref{subsec:causas-erros} \\
Completude do código gerado (Métricas 1.1 a 1.8) & 546 Métodos de Passo e 428 Métodos POM, nenhum sem implementação (100\%) & \ref{subsec:eficacia} \\
Cenários sem erro na 1ª execução (Métrica 1.9) & 12 de 27 (44,4\%) & \ref{subsec:qualidade} \\
\hline
\end{tabular}
\legend{Fonte: Autores (2026).}
\end{table}

\subsection{Eficácia: completude estrutural do código gerado (Métricas 1.1 a 1.8)}
\label{subsec:eficacia}

A eficácia avalia se a IA produz código estruturalmente completo a partir do cenário BDD, ou seja, quantos Métodos de Passo e Métodos POM foram gerados (Métricas 1.1 e 1.3), quantos permaneceram sem implementação (Métricas 1.2 e 1.4) e as respectivas frações (Métricas 1.5 a 1.8). A Tabela~\ref{tab:completude} apresenta os resultados.

\begin{table}[htb]
\centering
\caption{Completude do código gerado (Métricas 1.1 a 1.8)}
\label{tab:completude}
\begin{tabular}{cp{7.2cm}p{4.6cm}}
\hline
\textbf{Métrica} & \textbf{Descrição} & \textbf{Resultado (54 tentativas válidas)} \\
\hline
1.1 & Métodos de Passo gerados & 546 \\
1.2 & Métodos de Passo sem implementação & 0 \\
1.3 & Métodos POM gerados & 428 \\
1.4 & Métodos POM sem implementação & 0 \\
1.5 & Fração de Métodos de Passo com implementação & 100\% em todas as tentativas \\
1.6 & Fração de Métodos de Passo sem implementação & 0\% em todas as tentativas \\
1.7 & Fração de Métodos POM com implementação & 100\% em todas as tentativas \\
1.8 & Fração de Métodos POM sem implementação & 0\% em todas as tentativas \\
\hline
\end{tabular}
\legend{Fonte: Autores (2026).}
\end{table}

Foram gerados 546 Métodos de Passo e 428 Métodos POM, e nenhum método, em nenhuma das duas categorias, ficou sem implementação. Os totais somam todas as tentativas e, como o código é gerado novamente a cada tentativa, não representam métodos distintos. Por tentativa, a IA gerou de 5 a 15 Métodos de Passo (média de 10,1) e de 4 a 12 Métodos POM (média de 7,9). A completude foi de 100\% em todas as 54 tentativas, inclusive nas 27 que falharam na execução, e já havia sido de 100\% nas três etapas de calibração do \textit{prompt} de sistema.

A eficácia, portanto, não variou entre os cenários, e não há diferença a ser explicada por essa métrica. Isso não permite afirmar que a escrita do BDD seja irrelevante para a completude, pois a ausência de variação indica um efeito teto: nas condições do estudo (modelo, ferramenta e \textit{prompt} de sistema), a métrica ficou saturada e não discriminou cenários mais bem ou mais mal escritos. Além disso, ela mede apenas a completude estrutural: um método com corpo implementado conta como completo mesmo que o seletor utilizado esteja errado e o teste falhe, como ocorreu nas 27 execuções com falha. A influência da escrita do BDD aparece, por isso, na qualidade.

\subsection{Qualidade: erros sintáticos e de execução (Métrica 1.9)}
\label{subsec:qualidade}

Em \citeonline[p.~3222 e 3224]{Nettur2025}, o código gerado é integrado ao repositório sem modificações iniciais, os testes são executados e os erros registrados no \textit{log} do Cypress são contabilizados. Para reproduzir essa condição, a Métrica 1.9 foi avaliada na primeira execução de cada cenário, isto é, no código tal como foi gerado, antes de qualquer correção. Um cenário é contado como livre de erros se o resultado dessa execução foi \textbf{S}. As execuções seguintes, resultantes das correções feitas, não entram na métrica, mas sustentam a análise da distribuição e das causas dos erros (Seções~\ref{sec:analise-descritiva} e~\ref{sec:analises-complementares}).

Apenas 12 dos 27 cenários (44,4\%) executaram sem erros na primeira execução, e os outros 15 (55,6\%) tiveram ao menos um erro antes de funcionar, somando 27 execuções com erro (N) em 54. Como mostrou a Seção~\ref{sec:analise-descritiva}, a maior parte dos cenários (25 de 27, ou 92,6\%) teve no máximo dois erros, e o pior caso foi a US06-CA00, com cinco (Figura~\ref{fig:tentativas-por-us}). A Seção~\ref{sec:analises-complementares} mostrou de onde vêm esses erros: 63,0\% das falhas estão ligadas à forma como o elemento de interface foi referenciado no BDD, e 13 dos 15 cenários corrigidos foram resolvidos alterando a escrita do cenário.

Quatro ressalvas de método devem ser consideradas na leitura desse resultado:

\begin{enumerate}
    \item A primeira execução da base é a primeira com o \textit{prompt} de sistema definitivo, mas os cenários BDD já haviam sido escritos e ajustados na calibração. O valor de 44,4\% tende, por isso, a ser favorável em relação a uma primeira geração.
    \item O resultado de 100\% de cenários sem erro após as correções decorre do protocolo, no qual cada cenário foi ajustado até funcionar. Esse indicador apenas confirma o protocolo, pois toda execução bem-sucedida é, por definição, livre de erros. Por isso, ele não foi usado como medida de qualidade.
    \item A proporção de 50\% de execuções sem erro (27 de 54) equivale ao inverso da média de execuções por cenário (1 ÷ 2,0) e depende da política de correção adotada. Ela também serve apenas como contexto.
    \item Cada cenário teve uma única geração inicial, de modo que o resultado da primeira execução pode refletir também a variabilidade do modelo, e não apenas a escrita do BDD, como exemplificam a US16-CA00 e a US25-CA00 (Seções~\ref{subsec:correcoes} e~\ref{subsec:imprevisibilidade}).
\end{enumerate}

\subsection{Resposta à QE-1}
\label{subsec:resposta-qe1}

Em resposta à QE-1, não se observou influência da forma de escrita dos cenários BDD sobre a eficácia, pois a completude do código gerado foi de 100\% em todas as tentativas e, portanto, não houve variação a explicar. Em contrapartida, a forma de escrita esteve associada à qualidade: 63,0\% das falhas decorreram de como o elemento de interface foi referenciado no BDD, e 13 dos 15 cenários que precisaram de correção foram resolvidos alterando a escrita do cenário. A variabilidade do modelo, observada nos cenários US16-CA00 e US25-CA00, sugere que a escrita do BDD não explica sozinha a qualidade do código gerado. A Tabela~\ref{tab:resposta-qe1} resume a resposta por dimensão.

\begin{table}[htb]
\centering
\caption{Resposta à QE-1 por dimensão}
\label{tab:resposta-qe1}
\begin{tabular}{p{3.0cm}p{5.0cm}p{6.6cm}}
\hline
\textbf{Dimensão} & \textbf{Resultado} & \textbf{Leitura} \\
\hline
Eficácia (Métricas 1.1 a 1.8) & 100\% de completude em todas as 54 tentativas & Não se observou influência da escrita do BDD, pois não houve variação a explicar (efeito teto). \\
\hline
Qualidade (Métrica 1.9) & 12 de 27 cenários (44,4\%) sem erro na 1ª execução; 27 execuções com erro em 54 & A escrita do BDD esteve associada a ela: 63,0\% das falhas decorrem de como o elemento de interface foi referenciado, e 13 de 15 correções alteraram o BDD. \\
\hline
\end{tabular}
\legend{Fonte: Autores (2026).}
\end{table}

É importante destacar que, por se tratar de um estudo observacional de caso único, sem grupo de controle e com cenários escritos e corrigidos iterativamente, os resultados descrevem associação e um mecanismo plausível, e não um efeito causal quantificado. Além das ressalvas da Seção~\ref{subsec:qualidade}, o CAPJu não foi desenvolvido pensando em automação e carece de identificadores estáveis em elementos importantes da interface, o que pode ter ampliado o peso das falhas ligadas ao elemento de interface. O Capítulo~\ref{ch:conclusao} interpreta esses resultados à luz da literatura e responde à questão principal da pesquisa.

